\section{Temporal Learning}\label{sec:temporal-learning}

We formalize how a physical controller updates structure-class beliefs across sequential problem instances.

\subsection{Setup}

A controller faces a sequence of problem instances $I_1, I_2, \ldots$ drawn from an unknown distribution $D$ over problem structures.

After solving instance $I_k$ (or abstaining), the controller observes:
\begin{itemize}
\item Whether $I_k$ belonged to a tractable subcase
\item Which structure class $C_k$ was detected (if any)
\item The physical cost paid (substrate-dependent $\kappa$)
\end{itemize}

\begin{claim}{Bayesian structure detection}{th:bayesian-learning}
The controller maintains posterior $P(C \mid I_1, \ldots, I_k)$ over structure classes. Update rule: standard Bayesian update with likelihood $P(I_k \mid C)$.
\claimstamp{T\ref{th:bayesian-learning}}{RG}
\end{claim}

\leanmeta{\LH{CC4}}

Convergence: does $P(C \mid I_1, \ldots, I_k) \to \delta_{C*}$ as $k \to \infty$?

Answer: yes, if the structure classes are identifiable (distinct likelihood functions).

\begin{claim}{Identifiability implies convergence}{prop:identifiability-convergence}
If structure classes $C$ are identifiable (distinct likelihood functions), then $P(C \mid I_1, \ldots, I_k) \to \delta_{C*}$ as $k \to \infty$.
\claimstamp{P\ref{prop:identifiability-convergence}}{ID}
\end{claim}

\leanmeta{\LH{CC19}}

\subsection{PAC Learning Connection}

Sample complexity of structure detection:
How many instances must the controller solve before it can identify the structure class with probability $\geq 1 - \delta$ and error $\leq \varepsilon$?

\begin{claim}{Sample complexity conjecture}{conj:sample-complexity}
For the represented family (bounded actions, separable utility, tree structure), sample complexity is polynomial in $1/\varepsilon$ and $\log(1/\delta)$.
\claimstamp{J\ref{conj:sample-complexity}}{RG}
\end{claim}

This connects to PAC learning theory; the structure classes are the hypothesis space.

\subsection{Learning Reduces Abstention}

As the controller learns the structure class, it can apply regime-appropriate heuristics. The abstention frontier shrinks over time, not because the problem gets easier, but because the model of the problem gets more accurate.

Formally: let $A_k$ be the set of instances where integrity forces abstention at time $k$. Then:
\[
A_1 \supseteq A_2 \supseteq \ldots \supseteq A_\infty
\]
where $A_\infty =$ instances genuinely outside all tractable subcases.

\begin{claim}{Abstention reduction over time}{th:abstention-reduces}
$A_k$ is decreasing in $k$: $A_1 \supseteq A_2 \supseteq \ldots$.
\claimstamp{T\ref{th:abstention-reduces}}{RG}
\end{claim}

\leanmeta{\LH{CC68}}

The trajectory diverges based on $\kappa$ even though verdicts are $\kappa$-independent. Different substrates adapt at different rates, have different abstention thresholds, and make different solving/abstention trade-offs.

This connects to Paper 5: temporal learning produces costly signals. Each solved instance is a costly signal in Paper 5's sense: producing a correct solution has truth-dependent cost. Runtime/energy traces are costly signals; self-report is cheap talk.
